
\documentclass[11pt,a4paper]{article}

\usepackage[margin=2.4cm]{geometry}
\usepackage{amsmath, amssymb}
\usepackage{booktabs}
\usepackage{array}
\usepackage{enumitem}
\usepackage[hidelinks]{hyperref}

\title{Rangkuman Alur Trading-BEI: Dari Return sampai Output Backtest}
\author{}
\date{}

\begin{document}

\maketitle

\section*{Gambaran Umum}

Repo \texttt{trading-BEI} membangun strategi long-only untuk saham BEI/IDX.
Alurnya secara besar adalah:

\[
\text{data saham}
\rightarrow
\text{fitur}
\rightarrow
\text{model}
\rightarrow
\text{score}
\rightarrow
\text{portfolio}
\rightarrow
\text{return}
\rightarrow
\text{metrics}
\]

Strategi portfolio baru yang dibahas adalah:

\[
\boxed{
score > 0
\rightarrow
\text{ambil Top-N}
\rightarrow
\text{bobot pro-rata berdasarkan score}
}
\]

\section{Return}

Return adalah perubahan harga saham.

Jika harga naik dari 2.000 ke 2.100, maka:

\[
R_t = \frac{2100}{2000} - 1 = 0.05 = 5\%
\]

Di repo, return harian tidak dihitung langsung dari \texttt{close\_t / close\_{t-1}},
tetapi dari:

\[
\boxed{
R_t = \frac{close_t}{prev\_close_t} - 1
}
\]

\texttt{prev\_close} dipakai karena IDX menyediakan harga pembanding yang sudah
disesuaikan untuk corporate action seperti stock split.

\section{Log Return}

Return biasa disebut simple return:

\[
R_t = \frac{close_t}{prev\_close_t} - 1
\]

Di repo, return juga dinyatakan sebagai log return:

\[
\boxed{
r_t = \log\left(\frac{close_t}{prev\_close_t}\right)
}
\]

Contoh:

\[
r_t = \log\left(\frac{2100}{2000}\right)
= \log(1.05)
\approx 0.0488
\]

Kelebihan log return adalah bisa dijumlahkan antar periode.

Misal:

\[
\log\left(\frac{P_1}{P_0}\right)
+
\log\left(\frac{P_2}{P_1}\right)
=
\log\left(\frac{P_2}{P_0}\right)
\]

\section{Adjusted Log Price Path}

Adjusted log price path adalah akumulasi log return.

\[
\boxed{
A_t = \sum_{u \leq t} r_u
}
\]

Contoh:

\[
A_3 = r_1 + r_2 + r_3
\]

Gunanya adalah untuk menghitung return dari hari masuk ke hari keluar:

\[
\boxed{
y = A_{exit} - A_{entry}
}
\]

Jadi kalau:

\[
A_{entry}=0.01,
\qquad
A_{exit}=0.04
\]

maka:

\[
y = 0.04 - 0.01 = 0.03
\]

Artinya log return sekitar 3\%.

\section{Target \texorpdfstring{$y$}{y}}

Target adalah jawaban benar yang dipakai untuk melatih model.

Untuk saham \(i\) di hari \(t\), target ditulis:

\[
y_{t,i}
\]

Artinya:

\[
\text{return masa depan saham } i
\]

Dengan konfigurasi:

\[
execution\_lag = 1
\]

dan:

\[
horizon = 1
\]

model melihat data sampai hari \(t\), lalu baru bisa beli di hari \(t+1\).
Return target dihitung dari hari \(t+1\) sampai \(t+2\):

\[
\boxed{
y_{t,i} = A_{t+2,i} - A_{t+1,i}
}
\]

Jika \texttt{horizon = 5}, maka:

\[
\boxed{
y_{t,i} = A_{t+6,i} - A_{t+1,i}
}
\]

Intinya:

\[
\boxed{
y = \text{return masa depan yang benar-benar terjadi}
}
\]

\section{Input Model \texorpdfstring{$X$}{X}}

Model tidak melihat masa depan. Model hanya melihat data masa lalu sampai hari \(t\).

Jika:

\[
lookback = 60
\]

maka input model adalah:

\[
\boxed{
X_{t,i} = \text{fitur saham } i \text{ dari } t-59 \text{ sampai } t
}
\]

Contoh kalau hari ini hari ke-60:

\[
X_{60} = \text{hari 1 sampai 60}
\]

Jika \texttt{execution\_lag = 1} dan \texttt{horizon = 5}, maka targetnya:

\[
y_{60} = A_{66} - A_{61}
\]

Bukan:

\[
X = \text{hari 7 sampai 66}
\]

karena hari 61 sampai 66 belum terjadi pada hari ke-60.
Kalau model melihat data sampai hari 66, itu disebut \textit{look-ahead bias}.

\[
\boxed{
X = \text{masa lalu}
\qquad
y = \text{masa depan}
}
\]

\section{Model Menghasilkan Score}

Model menerima input:

\[
X_{t,i}
\]

lalu menghasilkan score:

\[
\boxed{
score_{t,i} = f_\theta(X_{t,i})
}
\]

Makna score:

\[
\boxed{
\text{score tinggi} = \text{saham lebih menarik menurut model}
}
\]

\[
\boxed{
\text{score rendah} = \text{saham kurang menarik menurut model}
}
\]

Model tidak langsung mengeluarkan keputusan beli atau jual.
Model hanya memberi angka score untuk setiap saham.

\section{Dari Score Menjadi Pilihan Saham}

Rule portfolio baru:

\[
\boxed{
score > 0
\rightarrow
\text{ambil Top-N}
\rightarrow
\text{bobot pro-rata}
}
\]

Misal:

\[
top\_n = 3
\]

\begin{center}
\begin{tabular}{lr}
\toprule
Saham & Score \\
\midrule
BBCA & 0.8 \\
TLKM & 0.3 \\
ASII & 0.1 \\
GOTO & -0.2 \\
BUKA & -0.5 \\
\bottomrule
\end{tabular}
\end{center}

Yang memenuhi \(score > 0\):

\[
BBCA,\ TLKM,\ ASII
\]

Total score positif:

\[
0.8 + 0.3 + 0.1 = 1.2
\]

Bobot:

\[
w_{BBCA} = \frac{0.8}{1.2}=66.7\%
\]

\[
w_{TLKM} = \frac{0.3}{1.2}=25\%
\]

\[
w_{ASII} = \frac{0.1}{1.2}=8.3\%
\]

Rumus umumnya:

\[
\boxed{
w_i =
\frac{score_i}
{\sum_{j \in P_t} score_j}
}
\]

dengan \(P_t\) adalah saham positif yang masuk Top-N.

Jika tidak ada saham dengan \(score > 0\), maka:

\[
\boxed{
100\% \text{ cash}
}
\]

\section{Dari Bobot Menjadi Portfolio Return}

Setelah saham dipilih, portfolio punya bobot:

\[
w_i
\]

Setiap saham punya return asli:

\[
R_i
\]

Portfolio return dihitung sebagai:

\[
\boxed{
R_p = \sum_i w_i R_i
}
\]

Contoh:

\begin{center}
\begin{tabular}{lrr}
\toprule
Saham & Bobot & Return Asli \\
\midrule
BBCA & 66.7\% & 2\% \\
TLKM & 25.0\% & -1\% \\
ASII & 8.3\% & 4\% \\
\bottomrule
\end{tabular}
\end{center}

Maka:

\[
R_p =
(0.667)(0.02)
+
(0.25)(-0.01)
+
(0.083)(0.04)
\]

\[
R_p = 0.01416 = 1.416\%
\]

\section{Transaction Cost}

Return sebelum biaya disebut gross return.

\[
R_{gross} = R_p
\]

Repo memakai biaya seperti:

\[
buy\_bps = 15
\qquad
sell\_bps = 25
\]

Artinya:

\[
buy = 0.15\%
\]

\[
sell = 0.25\%
\]

Jika buy turnover:

\[
0.8
\]

dan sell turnover:

\[
0.3
\]

maka cost:

\[
cost = 0.8(0.15\%) + 0.3(0.25\%)
\]

\[
cost = 0.12\% + 0.075\% = 0.195\%
\]

Net return:

\[
\boxed{
R_{net} = R_{gross} - cost
}
\]

Contoh:

\[
R_{net} = 1.416\% - 0.195\% = 1.221\%
\]

\section{Sharpe Ratio}

Sharpe mengukur return dibanding risiko.

Rumus sederhananya:

\[
Sharpe =
\frac{
mean(R_{net}) - r_f
}{
std(R_{net})
}
\]

Di repo, Sharpe di-annualized:

\[
\boxed{
Sharpe =
\frac{
mean(R_{net}) - r_f
}{
std(R_{net}) + \epsilon
}
\sqrt{252}
}
\]

Di sini \(r_f\) adalah risk-free rate \emph{per hari} (bukan per tahun), dan
Sharpe dihitung atas \emph{excess return} \(R_{net} - r_f\). Faktor \(\sqrt{252}\)
berlaku untuk rebalance harian; untuk cadence lain (mingguan/bulanan dengan
\texttt{rebalance\_every} = 5/21) faktornya \(\sqrt{252/\text{period\_days}}\),
yaitu akar dari jumlah periode per tahun.

Maknanya:

\begin{itemize}[leftmargin=*]
    \item return rata-rata tinggi membuat Sharpe naik;
    \item volatilitas tinggi membuat Sharpe turun;
    \item strategi yang stabil untung biasanya punya Sharpe lebih baik.
\end{itemize}

\section{Objective / Loss Model}

Tujuan model adalah membuat Sharpe setinggi mungkin.

Tetapi optimizer di machine learning biasanya meminimalkan loss.
Maka loss dibuat:

\[
\boxed{
loss = -Sharpe
}
\]

Jika Sharpe naik, loss turun.

\begin{center}
\begin{tabular}{rr}
\toprule
Sharpe & Loss \\
\midrule
0.5 & -0.5 \\
1.0 & -1.0 \\
2.0 & -2.0 \\
\bottomrule
\end{tabular}
\end{center}

Jadi model dilatih untuk:

\[
\boxed{
\text{memaksimalkan net Sharpe}
}
\]

\section{Backpropagation}

Setelah loss dihitung, model diperbaiki dengan backpropagation.

Model punya parameter:

\[
\theta
\]

Training melakukan:

\[
loss
\rightarrow
gradient
\rightarrow
update\ \theta
\]

Secara konsep:

\[
\boxed{
\theta_{baru}
=
\theta_{lama}
-
\text{learning rate} \times \text{gradient}
}
\]

Alur training:

\[
X
\rightarrow
model
\rightarrow
score
\rightarrow
portfolio
\rightarrow
net\ return
\rightarrow
Sharpe
\rightarrow
loss
\rightarrow
update\ model
\]

Intinya:

\[
\boxed{
\text{model belajar dari hasil portfolio, bukan hanya dari prediksi satu saham}
}
\]

\section{Training Portfolio vs Real Trading Rule}

Saat training, repo memakai softmax agar bisa dihitung gradient-nya.

\[
\boxed{
w_i =
\frac{e^{score_i/\tau}}
{\sum_j e^{score_j/\tau}}
}
\]

Softmax membuat score menjadi bobot yang smooth.

Saat backtest/trading rule, strategi memakai rule asli:

\[
\boxed{
score > 0
\rightarrow
\text{Top-N}
\rightarrow
\text{pro-rata}
}
\]

Jadi ada perbedaan:

\[
\text{training} = \text{softmax}
\]

\[
\text{backtest} = \text{positive Top-N pro-rata}
\]

Ini adalah mismatch, tetapi bukan bug.
Alasannya: Top-N dan filter \(score>0\) tidak smooth, sehingga sulit untuk backpropagation.

Mismatch ini dikurangi dengan:

\begin{itemize}[leftmargin=*]
    \item softmax temperature kecil, sehingga softmax lebih tajam;
    \item \texttt{allow\_cash = true}, sehingga cash punya logit tetap 0;
    \item model selection memakai rule trading asli.
\end{itemize}

\section{Arti \texorpdfstring{$score > 0$}{score > 0}}

Untuk model Sharpe-softmax, \(score > 0\) bukan berarti prediksi return positif.

Maknanya adalah:

\[
\boxed{
\text{score saham} > \text{score cash}
}
\]

Karena cash punya logit:

\[
s_{cash}=0
\]

Maka \(score > 0\) berarti saham dianggap lebih menarik daripada cash anchor.

Namun, jika model dilatih dengan objective regression/MSE, maka score bisa lebih dekat
maknanya ke prediksi return.

\section{Validation / Model Selection}

Training berjalan beberapa epoch.

Contoh:

\[
epoch = 1,2,3,\ldots,50
\]

Setiap epoch menghasilkan model berbeda.

Model yang dipakai bukan selalu epoch terakhir, tetapi model dengan validation Sharpe terbaik.

Contoh:

\begin{center}
\begin{tabular}{rr}
\toprule
Epoch & Validation Sharpe \\
\midrule
1 & 0.4 \\
2 & 0.8 \\
3 & 0.6 \\
\bottomrule
\end{tabular}
\end{center}

Model yang dipilih adalah epoch 2 karena Sharpe validation paling tinggi.

Catatan penting: validation Sharpe di sini dihitung dengan \emph{quick backtest}
yang frictionless --- memakai rule portfolio yang sama (mis. positive Top-N
pro-rata), tetapi \textbf{tanpa} ARA/ARB, spread, suspension, maupun delisting.
Aturan tradability penuh baru diterapkan di realistic backtest saat evaluasi
akhir. Jadi yang "sama" antara validation dan evaluasi adalah \emph{rule
portfolio}-nya, bukan keseluruhan simulasi pasar.

Intinya:

\[
\boxed{
\text{training belajar pakai softmax, tetapi model dipilih pakai rule trading asli}
}
\]

\section{Realistic Backtest}

Setelah model dipilih, model diuji di test set / out-of-sample.

Realistic backtest mensimulasikan kondisi trading yang lebih dekat dengan IDX.

Hal-hal yang diperhitungkan:

\begin{enumerate}[leftmargin=*]
    \item \textbf{Execution lag}:
    signal dari hari \(t\), beli di hari \(t+1\).
    \item \textbf{ARA}:
    tidak bisa beli jika tidak ada offer.
    \item \textbf{ARB}:
    tidak bisa jual jika tidak ada bid.
    \item \textbf{Suspension}:
    posisi tetap ada, tidak dihapus.
    \item \textbf{Delisting}:
    posisi bisa kena write-down.
    \item \textbf{Transaction cost dan spread}:
    biaya beli, biaya jual, dan bid-ask spread dihitung.
\end{enumerate}

Alurnya:

\[
score
\rightarrow
pilih\ saham
\rightarrow
cek\ bisa\ beli
\rightarrow
beli
\rightarrow
hitung\ return
\rightarrow
cek\ bisa\ jual
\rightarrow
hitung\ cost
\rightarrow
net\ return
\]

\section{Metrics Utama}

Setelah backtest, repo menghasilkan metrics.

\subsection{Annual Return}

\[
ann\_return
\]

Artinya return tahunan strategi.

Jika:

\[
ann\_return = 0.18
\]

maka strategi menghasilkan sekitar:

\[
18\% \text{ per tahun}
\]

\subsection{Annual Volatility}

\[
ann\_vol
\]

Artinya volatilitas tahunan.

Volatilitas tinggi berarti return lebih naik-turun.

\subsection{Sharpe}

Sharpe mengukur return dibanding risiko.

\[
\boxed{
\text{Sharpe tinggi} = \text{strategi lebih efisien}
}
\]

\subsection{Max Drawdown}

Max drawdown adalah penurunan terbesar dari puncak ke bawah.

Contoh equity:

\[
100 \rightarrow 130 \rightarrow 90
\]

Drawdown:

\[
\frac{90}{130} - 1 = -30.8\%
\]

Jika:

\[
max\_drawdown = -0.35
\]

artinya strategi pernah turun 35\% dari puncak.

\section{Metrics Lanjutan}

\subsection{Turnover}

Turnover mengukur seberapa banyak portfolio berubah.

\[
\boxed{
\text{turnover tinggi} = \text{sering beli/jual}
}
\]

Turnover tinggi bisa membuat cost tinggi.

\subsection{Cost}

Cost adalah biaya transaksi.

\[
cost = buy\ cost + sell\ cost + spread
\]

Jika cost tinggi:

\[
R_{net} \downarrow
\]

\subsection{Cash}

Cash adalah modal yang tidak dibelikan saham.

Dalam rule baru, cash muncul jika:

\[
\text{tidak ada saham dengan } score > 0
\]

atau saham yang memenuhi constraint tidak bisa dibeli.

\subsection{\texorpdfstring{$n\_stuck$}{n\_stuck}}

\[
n\_stuck
\]

adalah jumlah saham yang nyangkut.

Contoh: model ingin jual, tetapi saham terkena ARB atau suspend sehingga tidak bisa dijual.

\section{Compare vs IHSG}

Karena strategi ini long-only saham Indonesia, benchmark yang masuk akal adalah IHSG.

Pertanyaannya bukan hanya:

\[
\text{strategi untung atau tidak?}
\]

Tetapi:

\[
\boxed{
\text{strategi mengalahkan IHSG atau tidak?}
}
\]

Jika:

\[
ann\_return_{strategy} = 18\%
\]

dan:

\[
ann\_return_{IHSG} = 12\%
\]

maka strategi lebih tinggi:

\[
18\% - 12\% = 6\%
\]

dan:

\[
beats\_ihsg = True
\]

Kalau strategi return 8\% dan IHSG return 12\%, maka:

\[
beats\_ihsg = False
\]

\section{File Output}

Setelah backtest selesai, repo menyimpan hasil di:

\[
\texttt{results/<experiment\_name>/}
\]

File utama:

\begin{enumerate}[leftmargin=*]
    \item \texttt{metrics.json}
    \item \texttt{daily\_returns.csv}
    \item \texttt{test\_scores.csv}
    \item \texttt{equity\_vs\_ihsg.png}
\end{enumerate}

\subsection{\texttt{test\_scores.csv}}

Berisi score model untuk saham di periode test.

Contoh isi:

\begin{center}
\begin{tabular}{llrr}
\toprule
date & ticker & score & fwd\_return \\
\midrule
2025-01-01 & BBCA & 0.8 & 0.02 \\
2025-01-01 & TLKM & 0.3 & -0.01 \\
\bottomrule
\end{tabular}
\end{center}

File ini dipakai untuk menentukan portfolio:

\[
score > 0
\rightarrow
TopN
\rightarrow
pro\ rata
\]

\subsection{\texttt{daily\_returns.csv}}

Berisi hasil simulasi portfolio harian.

Contoh:

\begin{center}
\begin{tabular}{lrrrrr}
\toprule
date & gross & cost & net & turnover & cash \\
\midrule
2025-01-01 & 0.012 & 0.002 & 0.010 & 0.8 & 0.0 \\
\bottomrule
\end{tabular}
\end{center}

Hubungannya:

\[
\boxed{
net = gross - cost
}
\]

\subsection{\texttt{metrics.json}}

Berisi ringkasan performa akhir.

Contoh:

\[
ann\_return = 0.18
\]

\[
Sharpe = 1.2
\]

\[
max\_drawdown = -0.15
\]

\[
beats\_ihsg = True
\]

Ini file paling cepat untuk melihat apakah strategi bagus atau tidak.

\subsection{\texttt{equity\_vs\_ihsg.png}}

Grafik yang membandingkan:

\[
equity\ strategi
\]

dengan:

\[
IHSG
\]

Dari grafik ini bisa dilihat apakah strategi benar-benar mengalahkan market.

\section*{Ringkasan Akhir}

Alur lengkapnya:

\[
close,\ prev\_close
\rightarrow
return
\rightarrow
log\ return
\rightarrow
A_t
\rightarrow
X,y
\rightarrow
model
\rightarrow
score
\rightarrow
portfolio
\rightarrow
gross\ return
\rightarrow
cost
\rightarrow
net\ return
\rightarrow
Sharpe
\rightarrow
loss
\rightarrow
training
\rightarrow
backtest
\rightarrow
metrics
\]

Strategi baru:

\[
\boxed{
score > 0
\rightarrow
\text{Top-N}
\rightarrow
\text{pro-rata}
\rightarrow
\text{cash jika kosong}
}
\]

Makna penting:

\[
\boxed{
score > 0
\neq
\text{prediksi return positif}
}
\]

Untuk training Sharpe dengan \texttt{allow\_cash=true}, maknanya adalah:

\[
\boxed{
\text{score saham} > \text{score cash}
}
\]

\end{document}
